\section{Evaluation}
\label{sec:results}

The evaluation matrix contains 30 cells spanning five architectures, two
workloads, and three contention levels. Table~\ref{tab:agg} reports
aggregate means over the full matrix. Figure~\ref{fig:mainresults}
focuses on the captured workload, which is the sharper stress case for
the publication channel. Figure~\ref{fig:workloads} then shows how the
ranking changes, or fails to change, when the workload shifts.

\begin{figure*}[!t]
  \centering
  \includegraphics[width=\textwidth]{figures/paper_main_results.pdf}
  \caption{Main result matrix for the captured sparse-divergence
  workload as contention increases from 0 to 4 stressor CPUs on the
  x-axis. Each panel compares the same five publication architectures
  under the same workload and stressor setting. Panel (a) plots the
  fraction of attempted monitor reads identified as torn snapshots;
  values above zero indicate that the publication method can return a
  mixed logical record. Panel (b) plots CPU-side coherence traffic per
  publish from Ruby message counts, so larger values mean more pressure
  on the shared hierarchy. Panel (c) plots whole-workload simulated time
  to retire 5000 publications, which we use as a throughput-style proxy
  for publication cost rather than as a latency bound. Panel (d) plots
  read retries per run on a symlog scale, making visible the difference
  between designs that almost never reject a read and designs that rely
  heavily on reader retries to preserve correctness.}
  \label{fig:mainresults}
\end{figure*}

\begin{table*}[t]
\centering
\small
\begin{tabular}{@{}lrrrrrr@{}}
\toprule
Architecture & Torn (\%) & CPU msgs/pub & DMA msgs/pub & Sim time (ms) & Retries/run & Accepted reads (\%) \\
\midrule
unsync       & 4.26  & 61.05 & 0.000 & 0.864 & 0.0     & 99.64 \\
dma-naive    & 18.22 & 49.76 & 0.356 & 0.891 & 0.0     & 25.37 \\
dma-seqlock  & 0.00  & 28.89 & 0.193 & 0.482 & 14344.8 & 8.03 \\
seqlock      & 0.00  & 37.76 & 0.000 & 0.485 & 1543.2  & 97.85 \\
dblbuf       & 0.00  & 27.61 & 0.000 & 0.346 & 0.7     & 99.95 \\
\bottomrule
\end{tabular}
\caption{Aggregate results across the 30-cell matrix, averaged over
  workloads and contention levels. CPU-side coherence traffic excludes
  DMA-controller messages. Simulated time is the whole-workload time for
  5000 publications. ``Accepted reads'' is the fraction of read attempts
  the monitor accepted; for unsynchronized architectures it is
  misleadingly high because torn reads are counted as accepted.}
\label{tab:agg}
\end{table*}

\subsection{Correctness}

Figure~\ref{fig:mainresults}, panel (a), shows the central correctness
result. The two architectures without a publication primitive, \unsync
and \dmanaive, return torn snapshots in every measured condition.
Averaged over the full matrix, \unsync tears on 4.26\% of reads and
\dmanaive on 18.22\%. All three disciplined architectures, \seqlock,
\dblbuf, and \dmaseqlock, return zero torn reads in every cell.

\dmanaive is especially informative. Explicit transfer does not repair a
broken publication path. Without a handshake, it adds a second race on
the mirror side and tears more often than the direct-sharing baseline.
Once the same synchronization discipline is imposed on the mirror,
correctness is restored. The control variable is therefore the
publication primitive, not the transport by itself.

\subsection{Cost of Correct Designs}

Once correctness is enforced, the three disciplined architectures
separate on cost. Panels (b) through (d) of
Figure~\ref{fig:mainresults} make that split visible on the captured
workload, while Table~\ref{tab:agg} shows that the same ordering holds
on average over the full matrix.

\dblbuf is the strongest overall point in the design space. Averaged
over all 30 cells, it has the lowest CPU-side coherence traffic among
the correct designs, 27.61 messages per publish, and the shortest
whole-workload simulated time, 0.346 ms. Relative to \seqlock,
\dblbuf reduces CPU-side traffic by 27\% and shortens whole-workload
time by 29\%. Relative to \dmaseqlock, it shortens whole-workload time
by 28\% while avoiding the synchronized DMA mirror's extra transfer
path.

\seqlock is also correct, but it keeps the producer and monitor on one
active record and therefore pays noticeably more CPU-side traffic,
37.76 messages per publish on average, together with about 1543 retries
per run. \dmaseqlock shifts most of the reader's interaction away from
the producer's active line and therefore lands close to \dblbuf in
CPU-side traffic, 28.89 messages per publish, but it replaces that
benefit with mirror refreshes and a very heavy retry burden. Averaged
over the full matrix, \dmaseqlock incurs about 14{,}345 retries per run,
more than 9$\times$ the retry count of \seqlock. The synchronized DMA
mirror is therefore correct and competitive, but it is not the cheapest
correct mechanism in this benchmark.

That comparison should be read carefully. \dmaseqlock is not a
``coherence disabled'' baseline. The mirror still lives in the same
coherent substrate. What it changes is where correctness work is paid:
the direct-sharing designs synchronize at the producer-owned record,
whereas the DMA design inserts a periodic copy stage and then validates
the mirror. The cost difference therefore reflects the architectural
price of turning one publication problem into two coupled stages,
producer to mirror and mirror to monitor.

\subsection{Read-Side Efficiency}

Retry count is only one side of reader cost. Table~\ref{tab:agg} also
reports the fraction of read attempts returned as data rather than as a
retry. On that metric the difference between the correct designs becomes
sharper. \dblbuf accepts 99.95\% of read attempts, \seqlock accepts
97.85\%, and \dmaseqlock accepts only 8.03\%.

The explanation is architectural rather than accidental. In \seqlock the
reader races with producer updates on one active record, so retries
occur only when the read overlaps that short update window. In
\dmaseqlock the reader polls a mirror that is refreshed periodically
regardless of the reader's progress, so every refresh creates another
opportunity to observe an in-progress publication and retry. That is why
\dmaseqlock can be logically correct yet operationally expensive for a
small supervisory core. If monitor CPU budget matters, consumer
efficiency is part of the channel design, not just a secondary concern.

\subsection{Workload Sensitivity}

\begin{figure}[t]
  \centering
  \includegraphics[width=0.98\columnwidth]{figures/paper_workload_sensitivity.pdf}
  \caption{Effect of workload shape on CPU-side traffic for the correct
  designs, averaged over stressor counts. The sustained-bias workload
  increases traffic for the direct coherent protocols, whereas the
  synchronized DMA mirror stays nearly flat.}
  \label{fig:workloads}
\end{figure}

Figure~\ref{fig:workloads} adds a second result. To keep the comparison
focused, it averages over stressor counts and shows only the correct
architectures. Workload shape changes cost more than correctness. Moving
from the captured sparse-divergence workload to the synthesized
sustained-bias workload increases average CPU-side coherence traffic by
37\% for \dblbuf and 36\% for \seqlock, which is the main effect shown
in the figure.
The whole-workload simulated time also rises, by 14\% for \dblbuf and
12\% for \seqlock. In contrast, the synchronized DMA mirror is almost
flat in CPU-side traffic: its average CPU-side message count changes by
less than 1\% across the two workloads.

That contrast is consistent with the transport choice. The two direct
coherent protocols expose the monitor more directly to producer activity
on the shared publication region, whereas \dmaseqlock makes the monitor
interact primarily with the mirror and leaves the producer-facing work
to the DMA engine. Even so, the global ranking does not change.
\dblbuf remains the lowest-cost correct design, \seqlock remains the
more expensive direct-sharing alternative, and \dmaseqlock remains a
correct but retry-heavy explicit-transfer design.

\subsection{Contention Sensitivity and Design Implications}

Increasing the number of stressor CPUs raises traffic and whole-workload
time for every architecture, as Figure~\ref{fig:mainresults} panels (b)
and (c) show. What does not change is the qualitative ordering. The
three disciplined architectures remain zero-torn, the two undisciplined
architectures remain incorrect, and \dblbuf remains the fastest design
in every measured cell.

The absolute growth under contention is also informative. On the
captured workload, moving from zero to four stressor CPUs raises
\dblbuf's CPU-side traffic from 12.84 to 32.70 messages per publish and
raises \seqlock's from 18.49 to 44.05. \dmaseqlock starts from the
lowest CPU-side traffic at zero stressors, 9.80, but climbs to 48.29 by
four stressors because the mirror path and the shared hierarchy are both
being exercised under interference. Direct coherent publication is
therefore not automatically lower-traffic than explicit transfer; the
interaction between transport, primitive, and interference still
matters.

The ranking nevertheless has a clear structural explanation. \dblbuf
keeps the writer and reader mostly off each other's active payload by
letting the producer fill an inactive slot and then publish a single
generation word. \seqlock keeps both parties concentrated on one hot
record, so validation and payload movement contend on the same line.
\dmaseqlock removes the reader from the producer's active record, but it
reintroduces overlap at the mirror through the periodic pull cadence.
What wins in this benchmark is therefore not ``direct sharing'' in the
abstract, but a specific combination of direct sharing with a primitive
that minimizes both invalidation pressure and read-side speculation.

Taken together, the results point to a simple design implication for
\CPS monitor channels. If the monitor is meant to consume the latest
producer record directly over a coherent fabric, the record should be
published with an explicit primitive and preferably with a
generation-based double buffer. If the system instead prefers an
explicit-transfer design for isolation reasons, that path still needs a
record-level handshake; transport alone does not provide correctness.
